%%%%%%%%%%%%%%%%%%%%%%% file template.tex %%%%%%%%%%%%%%%%%%%%%%%%%
%
% This is a template file for Web of Conferences Journal
%
% Copy it to a new file with a new name and use it as the basis
% for your article
%
%%%%%%%%%%%%%%%%%%%%%%%%%% EDP Science %%%%%%%%%%%%%%%%%%%%%%%%%%%%
%
%%%\documentclass[option comma separated list]{webofc}
%%%Three important options:
%%% "epj" for EPJ Web of Conferences Journal
%%% "bio" for BIO Web of Conferences Journal
%%% "mat" for MATEC Web of Conferences Journal
%%% "itm" for ITM Web of Conferences Journal
%%% "e3s" for E3S Web of Conferences Journal
%%% "shs" for SHS Web of Conferences Journal
%%% "twocolumn" for typesetting an article in two columns format (default one column)
%
\documentclass[aas_macros]{webofc}
\usepackage[varg]{txfonts}   % Web of Conferences font
\usepackage{listings}
\usepackage{hyperref}
\lstset
{
    language=[LaTeX]TeX,
    breaklines=true,
    basicstyle=\tt\small
}
%
% Put here some packages required or/and some personnal commands

\newcommand{\todo}[1]{{\bf #1}}
\newcommand{\bigO}{\mathcal{O}}
\newcommand{\savg}[1]{\left<#1\right>}
\newcommand{\svar}{{\rm Var}}
\newcommand{\scov}{{\rm Cov}}
\newcommand{\Mshft}{\mathbf{M}_{\theta_2}}
\newcommand{\dMshft}{\partial\Mshft}
\newcommand{\dA}{\partial A}
\newcommand{\dB}{\partial B}
%
%
\begin{document}
%
\title{A fast template periodogram}
%
% subtitle is optionnal
%
%%%\subtitle{Do you have a subtitle?\\ If so, write it here}

\author{\firstname{John} \lastname{Hoffman}\inst{1}\fnsep\thanks{\href{mailto:jah5@princeton.edu}{\tt jah5@princeton.edu}}}
\institute{Department of Astrophysical Sciences, Princeton University, Princeton NJ 08540}

\abstract{%
  This proceedings presents a novel, non-linear extension to the Lomb-Scargle periodogram that allows periodograms to be generated
  for arbitrary signal shapes. Such periodograms are already known as "template periodograms" or "periodic matched filters," but 
  current implementations are computationally inefficient. The "fast template periodogram" presented here improves existing techniques
  by a factor of $\sim$a few for small test cases ($\bigO(10)$ observations), and over three orders of magnitude for lightcurves containing
  $\bigO(10^4)$ observations. The fast template periodogram scales asymptotically as $\bigO(HN_f\log HN_f + H^4N_f)$, where
  $H$ denotes the number of harmonics required to adequately approximate the template and $N_f$ is the number of trial
  frequencies. Existing implementations scale as $\bigO(N_{\rm obs}N_f)$, where $N_{\rm obs}$ is the number of
  observations in the lightcurve. An open source Python implementation is available on GitHub.
}
%
\maketitle
%


\section{Introduction}\label{sec:intro}
Template fitting, or periodic matched filtering, is a powerful technique for identifying, classifying, and characterizing
periodic variability in noisy time series data. However, template fitting requires several orders of magnitude more computational
resources than running a Lomb Scargle periodogram \cite{Lomb_1976,Scargle_1982,Barning_1963,Vanicek_1971} on the same lightcurve, 
and the computational requirements scale as $N_{\rm obs}N_f$, where $N_{\rm obs}$ is the number of observations in the 
lightcurve and $N_f$ is the number of trial frequencies.

For example, \cite{Sesar_et_al_2016} used template fitting to model RR Lyrae in the Pan-STARRS DR1 \citep{PanSTARRS} and found that, 
for a given completeness, template fitting produced purer samples of RR Lyrae than Lomb-Scargle or multi-term extensions \citep{Palmer_2009}.
However, they found that the computational resources needed to produce template fits (approximately 30 minutes per source per CPU), forced
them to 





%\begin{figure*}
%\centering
% Use the relevant command for your figure-insertion program
% to insert the figure file. See example above.
% If not, use
%\includegraphics[width=10cm,clip]{figure-epjwc.pdf}
%\caption{Please write your figure caption here  (standard formatting for the figure and caption).}
%\label{fig:fig-1}       % Give a unique label
%\end{figure*}

%\begin{figure}
% Use the relevant command for your figure-insertion program
% to insert the figure file.
%\centering
%\sidecaption
%\includegraphics[width=6cm,clip]{figure-epjwc.pdf}
%\caption{Please write your figure caption here (figure and caption side-by-side).}
%\label{fig:fig-2}       % Give a unique label
%\end{figure}

%\begin{table}
%\centering
%\caption{Please write your table header here.}
% For LaTeX tables you can use
%\begin{tabular}{lll}
%\hline
%first & second & third  \\\hline
%number & number & number \\
%number & number & number \\\hline
%\end{tabular}
%\label{tab:tab-1}       % Give a unique label
%\end{table}
%

%\subsection{Equations}\label{sec:eqs}
%For typesetting equations, you can use the usual commands in \LaTeX. Please keep in mind that equations should read as part of the flow of the text and be punctuated as such. For instance, if you write 



\begin{acknowledgement} 
\noindent\vskip 0.2cm
\noindent {\em Acknowledgments}: 
\end{acknowledgement}

\bibliography{refs}
\end{document}

% end of file template.tex

